% Options for packages loaded elsewhere
\PassOptionsToPackage{unicode}{hyperref}
\PassOptionsToPackage{hyphens}{url}
\documentclass[
]{article}
\usepackage{xcolor}
\usepackage[margin=1in]{geometry}
\usepackage{amsmath,amssymb}
\setcounter{secnumdepth}{-\maxdimen} % remove section numbering
\usepackage{iftex}
\ifPDFTeX
  \usepackage[T1]{fontenc}
  \usepackage[utf8]{inputenc}
  \usepackage{textcomp} % provide euro and other symbols
\else % if luatex or xetex
  \usepackage{unicode-math} % this also loads fontspec
  \defaultfontfeatures{Scale=MatchLowercase}
  \defaultfontfeatures[\rmfamily]{Ligatures=TeX,Scale=1}
\fi
\usepackage{lmodern}
\ifPDFTeX\else
  % xetex/luatex font selection
\fi
% Use upquote if available, for straight quotes in verbatim environments
\IfFileExists{upquote.sty}{\usepackage{upquote}}{}
\IfFileExists{microtype.sty}{% use microtype if available
  \usepackage[]{microtype}
  \UseMicrotypeSet[protrusion]{basicmath} % disable protrusion for tt fonts
}{}
\makeatletter
\@ifundefined{KOMAClassName}{% if non-KOMA class
  \IfFileExists{parskip.sty}{%
    \usepackage{parskip}
  }{% else
    \setlength{\parindent}{0pt}
    \setlength{\parskip}{6pt plus 2pt minus 1pt}}
}{% if KOMA class
  \KOMAoptions{parskip=half}}
\makeatother
\usepackage{color}
\usepackage{fancyvrb}
\newcommand{\VerbBar}{|}
\newcommand{\VERB}{\Verb[commandchars=\\\{\}]}
\DefineVerbatimEnvironment{Highlighting}{Verbatim}{commandchars=\\\{\}}
% Add ',fontsize=\small' for more characters per line
\usepackage{framed}
\definecolor{shadecolor}{RGB}{248,248,248}
\newenvironment{Shaded}{\begin{snugshade}}{\end{snugshade}}
\newcommand{\AlertTok}[1]{\textcolor[rgb]{0.94,0.16,0.16}{#1}}
\newcommand{\AnnotationTok}[1]{\textcolor[rgb]{0.56,0.35,0.01}{\textbf{\textit{#1}}}}
\newcommand{\AttributeTok}[1]{\textcolor[rgb]{0.13,0.29,0.53}{#1}}
\newcommand{\BaseNTok}[1]{\textcolor[rgb]{0.00,0.00,0.81}{#1}}
\newcommand{\BuiltInTok}[1]{#1}
\newcommand{\CharTok}[1]{\textcolor[rgb]{0.31,0.60,0.02}{#1}}
\newcommand{\CommentTok}[1]{\textcolor[rgb]{0.56,0.35,0.01}{\textit{#1}}}
\newcommand{\CommentVarTok}[1]{\textcolor[rgb]{0.56,0.35,0.01}{\textbf{\textit{#1}}}}
\newcommand{\ConstantTok}[1]{\textcolor[rgb]{0.56,0.35,0.01}{#1}}
\newcommand{\ControlFlowTok}[1]{\textcolor[rgb]{0.13,0.29,0.53}{\textbf{#1}}}
\newcommand{\DataTypeTok}[1]{\textcolor[rgb]{0.13,0.29,0.53}{#1}}
\newcommand{\DecValTok}[1]{\textcolor[rgb]{0.00,0.00,0.81}{#1}}
\newcommand{\DocumentationTok}[1]{\textcolor[rgb]{0.56,0.35,0.01}{\textbf{\textit{#1}}}}
\newcommand{\ErrorTok}[1]{\textcolor[rgb]{0.64,0.00,0.00}{\textbf{#1}}}
\newcommand{\ExtensionTok}[1]{#1}
\newcommand{\FloatTok}[1]{\textcolor[rgb]{0.00,0.00,0.81}{#1}}
\newcommand{\FunctionTok}[1]{\textcolor[rgb]{0.13,0.29,0.53}{\textbf{#1}}}
\newcommand{\ImportTok}[1]{#1}
\newcommand{\InformationTok}[1]{\textcolor[rgb]{0.56,0.35,0.01}{\textbf{\textit{#1}}}}
\newcommand{\KeywordTok}[1]{\textcolor[rgb]{0.13,0.29,0.53}{\textbf{#1}}}
\newcommand{\NormalTok}[1]{#1}
\newcommand{\OperatorTok}[1]{\textcolor[rgb]{0.81,0.36,0.00}{\textbf{#1}}}
\newcommand{\OtherTok}[1]{\textcolor[rgb]{0.56,0.35,0.01}{#1}}
\newcommand{\PreprocessorTok}[1]{\textcolor[rgb]{0.56,0.35,0.01}{\textit{#1}}}
\newcommand{\RegionMarkerTok}[1]{#1}
\newcommand{\SpecialCharTok}[1]{\textcolor[rgb]{0.81,0.36,0.00}{\textbf{#1}}}
\newcommand{\SpecialStringTok}[1]{\textcolor[rgb]{0.31,0.60,0.02}{#1}}
\newcommand{\StringTok}[1]{\textcolor[rgb]{0.31,0.60,0.02}{#1}}
\newcommand{\VariableTok}[1]{\textcolor[rgb]{0.00,0.00,0.00}{#1}}
\newcommand{\VerbatimStringTok}[1]{\textcolor[rgb]{0.31,0.60,0.02}{#1}}
\newcommand{\WarningTok}[1]{\textcolor[rgb]{0.56,0.35,0.01}{\textbf{\textit{#1}}}}
\usepackage{graphicx}
\makeatletter
\newsavebox\pandoc@box
\newcommand*\pandocbounded[1]{% scales image to fit in text height/width
  \sbox\pandoc@box{#1}%
  \Gscale@div\@tempa{\textheight}{\dimexpr\ht\pandoc@box+\dp\pandoc@box\relax}%
  \Gscale@div\@tempb{\linewidth}{\wd\pandoc@box}%
  \ifdim\@tempb\p@<\@tempa\p@\let\@tempa\@tempb\fi% select the smaller of both
  \ifdim\@tempa\p@<\p@\scalebox{\@tempa}{\usebox\pandoc@box}%
  \else\usebox{\pandoc@box}%
  \fi%
}
% Set default figure placement to htbp
\def\fps@figure{htbp}
\makeatother
\setlength{\emergencystretch}{3em} % prevent overfull lines
\providecommand{\tightlist}{%
  \setlength{\itemsep}{0pt}\setlength{\parskip}{0pt}}
\usepackage{bookmark}
\IfFileExists{xurl.sty}{\usepackage{xurl}}{} % add URL line breaks if available
\urlstyle{same}
\hypersetup{
  pdftitle={Multivariate Data Analysis Assignment 3},
  hidelinks,
  pdfcreator={LaTeX via pandoc}}

\title{Multivariate Data Analysis Assignment 3}
\author{}
\date{\vspace{-2.5em}}

\begin{document}
\maketitle

\begin{Shaded}
\begin{Highlighting}[]
\FunctionTok{source}\NormalTok{(}\AttributeTok{file =} \StringTok{"./R\_SCRIPTS/Functions.R"}\NormalTok{)}

\FunctionTok{library}\NormalTok{(tidyverse)}
\end{Highlighting}
\end{Shaded}

\begin{verbatim}
## -- Attaching core tidyverse packages ------------------------ tidyverse 2.0.0 --
## v dplyr     1.1.4     v readr     2.1.5
## v forcats   1.0.1     v stringr   1.5.2
## v ggplot2   4.0.0     v tibble    3.3.0
## v lubridate 1.9.4     v tidyr     1.3.1
## v purrr     1.1.0     
## -- Conflicts ------------------------------------------ tidyverse_conflicts() --
## x dplyr::filter() masks stats::filter()
## x dplyr::lag()    masks stats::lag()
## i Use the conflicted package (<http://conflicted.r-lib.org/>) to force all conflicts to become errors
\end{verbatim}

\begin{Shaded}
\begin{Highlighting}[]
\NormalTok{D }\OtherTok{\textless{}{-}} \FunctionTok{load\_data}\NormalTok{()}
\end{Highlighting}
\end{Shaded}

\subsection{Introduction}\label{introduction}

\subsection{Part 1: Exploratory Data
Analysis}\label{part-1-exploratory-data-analysis}

The data that we analyse for this report is mRNA gene expression data,
from 60 samples and across 2690 genes from a hypothetical (simulated)
study.

The aims of this simulated study are to analyse the effects of two
treatments on a cell culture. The 60 samples are split across two
covariates: time (6 or 12 hours) and treatment (lisinopril, amlodipine,
or control).

\begin{Shaded}
\begin{Highlighting}[]
\CommentTok{\# HISTOGRAM OF LOG GENE EXPRESSION IN A SAMPLE / FOR A PARTICULAR GENE}
\FunctionTok{ggplot}\NormalTok{() }\SpecialCharTok{+}
  \FunctionTok{geom\_histogram}\NormalTok{(}\AttributeTok{mapping =} \FunctionTok{aes}\NormalTok{(}\AttributeTok{x =}\NormalTok{ D}\SpecialCharTok{$}\NormalTok{data[}\DecValTok{1}\NormalTok{,]), }\AttributeTok{bins =} \DecValTok{40}\NormalTok{, }\AttributeTok{fill =} \StringTok{"steelblue"}\NormalTok{, }
                 \AttributeTok{colour =} \StringTok{"white"}\NormalTok{) }\SpecialCharTok{+}
  \FunctionTok{theme\_bw}\NormalTok{() }\SpecialCharTok{+}
  \FunctionTok{labs}\NormalTok{(}\AttributeTok{x =} \StringTok{"log Gene Expression"}\NormalTok{,}
       \AttributeTok{y =} \StringTok{"Frequency"}\NormalTok{,}
       \AttributeTok{title =} \StringTok{"Sa6C1: log Gene Expression Histogram"}\NormalTok{)}
\end{Highlighting}
\end{Shaded}

\begin{figure}
\centering
\pandocbounded{\includegraphics[keepaspectratio]{Report_files/figure-latex/fig-histogram-1.pdf}}
\caption{Histogram Displaying Logarithm of Gene Expression}
\end{figure}

\begin{Shaded}
\begin{Highlighting}[]
\CommentTok{\# BOXPLOT OF LOG GENE EXPRESSION IN A SAMPLE / FOR A PARTICULAR GENE}
\NormalTok{selection }\OtherTok{=} \FunctionTok{c}\NormalTok{(}\StringTok{"sample"}\NormalTok{, }\StringTok{"times"}\NormalTok{, }\StringTok{"treatments"}\NormalTok{)}

\NormalTok{long\_data }\OtherTok{\textless{}{-}} \FunctionTok{as.data.frame}\NormalTok{(D}\SpecialCharTok{$}\NormalTok{data) }\SpecialCharTok{|\textgreater{}} 
  \FunctionTok{add\_column}\NormalTok{(}\StringTok{"times"} \OtherTok{=}\NormalTok{ D}\SpecialCharTok{$}\NormalTok{times, }
             \StringTok{"treatments"} \OtherTok{=}\NormalTok{ D}\SpecialCharTok{$}\NormalTok{treatments) }\SpecialCharTok{|\textgreater{}} 
  \FunctionTok{rownames\_to\_column}\NormalTok{(}\AttributeTok{var =} \StringTok{"sample"}\NormalTok{) }\SpecialCharTok{|\textgreater{}} 
  \FunctionTok{pivot\_longer}\NormalTok{(}\AttributeTok{cols =} \SpecialCharTok{!}\FunctionTok{any\_of}\NormalTok{((selection)),}
               \AttributeTok{names\_to =} \StringTok{"gene"}\NormalTok{,}
               \AttributeTok{values\_to =} \StringTok{"expression"}\NormalTok{)}

\FunctionTok{ggplot}\NormalTok{(}\AttributeTok{data =}\NormalTok{ long\_data, }\FunctionTok{aes}\NormalTok{(}\AttributeTok{x =}\NormalTok{ treatments, }\AttributeTok{y =}\NormalTok{ expression, }\AttributeTok{fill =}\NormalTok{ treatments,}
                             \AttributeTok{group =}\NormalTok{ sample)) }\SpecialCharTok{+}
  \FunctionTok{geom\_boxplot}\NormalTok{(}\AttributeTok{outlier.size =} \FloatTok{0.4}\NormalTok{, }\AttributeTok{outlier.alpha =} \FloatTok{0.4}\NormalTok{) }\SpecialCharTok{+}
  \FunctionTok{facet\_wrap}\NormalTok{(}\SpecialCharTok{\textasciitilde{}}\NormalTok{times, }\AttributeTok{labeller =}\NormalTok{ label\_both) }\SpecialCharTok{+}
  \FunctionTok{scale\_fill\_manual}\NormalTok{(}\AttributeTok{values =}\NormalTok{ treatment\_colours, }\AttributeTok{guide =} \StringTok{"none"}\NormalTok{) }\SpecialCharTok{+}
  \FunctionTok{theme\_bw}\NormalTok{()}
\end{Highlighting}
\end{Shaded}

\begin{figure}
\centering
\pandocbounded{\includegraphics[keepaspectratio]{Report_files/figure-latex/fig-boxplot-1.pdf}}
\caption{Boxplot Displaying Distribution of Log Gene Expressions, Split
by Covariates}
\end{figure}

Figures @ref(fig:fig-histogram) and @ref(fig:fig-boxplot) display the
distributions of gene expression over our 60 samples. Note the logarithm
of the gene expression is given instead of the raw expression value, as
the data is given to us in log-normalised form. Then, when plotting the
data in our histogram, we see a distribution that resembles a Gaussian
one. This we expect.

We may compare differences between covariates in Figure
@ref(fig:fig-boxplot).

\begin{Shaded}
\begin{Highlighting}[]
\CommentTok{\# PLOT A TIME COURSE FOR A GENE IN DIFFERENT TREATMENTS}

\NormalTok{time\_data }\OtherTok{\textless{}{-}} \FunctionTok{as.data.frame}\NormalTok{(D}\SpecialCharTok{$}\NormalTok{data[, }\DecValTok{1}\NormalTok{]) }\SpecialCharTok{|\textgreater{}} 
  \FunctionTok{rename}\NormalTok{(}\AttributeTok{expression =} \DecValTok{1}\NormalTok{) }\SpecialCharTok{|\textgreater{}} 
  \FunctionTok{add\_column}\NormalTok{(}\StringTok{"times"} \OtherTok{=}\NormalTok{ D}\SpecialCharTok{$}\NormalTok{times,}
             \StringTok{"treatments"} \OtherTok{=}\NormalTok{ D}\SpecialCharTok{$}\NormalTok{treatments)}

\FunctionTok{ggplot}\NormalTok{(}\AttributeTok{data =}\NormalTok{ time\_data) }\SpecialCharTok{+}
  \FunctionTok{geom\_jitter}\NormalTok{(}\AttributeTok{mapping =} \FunctionTok{aes}\NormalTok{(}\AttributeTok{x =}\NormalTok{ times, }\AttributeTok{y =}\NormalTok{ expression, }\AttributeTok{colour =}\NormalTok{ treatments), }\AttributeTok{width =} \FloatTok{0.1}\NormalTok{) }\SpecialCharTok{+}
  \FunctionTok{geom\_line}\NormalTok{(}\AttributeTok{mapping =} \FunctionTok{aes}\NormalTok{(}\AttributeTok{x =}\NormalTok{ times, }\AttributeTok{y =}\NormalTok{ expression, }\AttributeTok{colour =}\NormalTok{ treatments,}
                          \AttributeTok{group =}\NormalTok{ treatments), }\AttributeTok{linewidth =} \FloatTok{0.8}\NormalTok{) }\SpecialCharTok{+}
  \FunctionTok{scale\_colour\_manual}\NormalTok{(}\AttributeTok{values =}\NormalTok{ treatment\_colours) }\SpecialCharTok{+}
  \FunctionTok{theme\_bw}\NormalTok{()}
\end{Highlighting}
\end{Shaded}

\pandocbounded{\includegraphics[keepaspectratio]{Report_files/figure-latex/Time Course-1.pdf}}

\subsection{Part 2: Prinicpal Components
Analysis}\label{part-2-prinicpal-components-analysis}

\subsection{Part 3: Cluster Analysis}\label{part-3-cluster-analysis}

\subsection{Part 4: Discriminant
Analysis}\label{part-4-discriminant-analysis}

\subsection{Part 5: Extension to Graph-Based
Clustering}\label{part-5-extension-to-graph-based-clustering}

\subsection{Bibliography}\label{bibliography}

\end{document}
